\documentclass[11pt]{article}
\usepackage[a4paper,margin=1in]{geometry}
\usepackage{amsmath,amssymb}
\usepackage{enumitem}
\usepackage{xcolor}

\definecolor{revcolor}{rgb}{0.0,0.3,0.6}
\newcommand{\rev}[1]{\textcolor{revcolor}{\emph{#1}}}

\title{Response to Reviewer Comments\\[4pt]
\large ``An Analytical First-Order Mixed-Zonal Theory for GEqOE''}
\author{Jose Manuel Montilla Garc\'ia}
\date{March 2026}

\begin{document}
\maketitle

We thank the referee for the thorough and constructive review. The
manuscript was submitted in internal-note form, and the revision addresses
all concerns raised. Below we provide a point-by-point response to each
finding.

\bigskip

%% ====================================================================
\section*{R-1: Harmonic-order contradiction ($|m| \le n$ vs $|m| \le n-2$)}

\rev{Reviewer: The paper states three different bounds for the maximum
harmonic order, which is confusing.}

\textbf{Response:} We agree. The preliminary bound $\mathcal{M}_n = \{m \le n\}$
has been immediately followed by a remark noting that the sharp bound
$|m| \le n-2$ applies (with forward reference to the proof). The total-model
summation bounds now include a remark that terms with $m > N-2$ have
identically zero coefficients. A new label \verb|\label{sec:harmonic_bound}|
has been added to the maximum-harmonic-order subsection for cross-referencing.

\textbf{Changes:} Eq.~after (Un\_fourier): added sharp-bound remark.
Section~10.2: added \verb|\label{sec:harmonic_bound}|.
Section~10.3 preamble: added remark on zero coefficients beyond $N-2$.

%% ====================================================================
\section*{R-2: Error-order inconsistency ($O(\varepsilon)$ vs $O(\varepsilon^2)$)}

\rev{Reviewer: The paper uses ``$O(\varepsilon^2)$'' and ``$O(J^2)$''
inconsistently across three separate validation contexts.}

\textbf{Response:} We agree. A three-part error taxonomy has been introduced
at the beginning of Section~11 (Validation), clearly distinguishing:
\begin{enumerate}[nosep]
\item Initialization error ($O(\varepsilon^2)$),
\item Short-period reconstruction error ($O(\varepsilon)$ amplitude),
\item Secular drift error ($O(\varepsilon^2)$ per orbit, linear accumulation).
\end{enumerate}

The slow-flow parity gap language has been corrected to distinguish absolute
and relative metrics. The scaling test opening sentence has been rewritten
to avoid the misleading ``expected $O(\varepsilon^2)$ residual'' phrasing;
it now states that the slope identifies the dominant error source.

\textbf{Changes:} Section~11: new error taxonomy (3-item enumerated list).
Section~10.5 (slow-flow parity): clarified relative vs absolute gap.
Section~11.4 (scaling test): rewritten opening paragraph.
Section~13 (conclusions): language consistent with corrected taxonomy.

%% ====================================================================
\section*{R-3: Fast-phase reconstruction ambiguity}

\rev{Reviewer: Steps 1 and 3 conflate $\bm{u}_1$ (slow state) and $v_1$
(fast phase). It is unclear when Kepler inversion is needed.}

\textbf{Response:} Agreed. Steps 1--3 of the operational procedure have been
rewritten to use both $\bm{u}_1$ and $v_1$ explicitly, with separate
equations for $\bar{\bm{z}}$ and $\bar K$. Step~2 now states that
$\bar M$ is propagated alongside the slow state and $\bar K$ is recovered
from $\bar M$ by solving the generalized Kepler equation. Step~3 clarifies
that the osculating $K$ is obtained directly via $v_1$ without an additional
Kepler inversion.

\textbf{Changes:} Section~9.2, Steps 1--3: rewritten with explicit
$\bm{z}$/$K$ notation. Section~11.2 (reconstruction procedure): step~4
clarified.

%% ====================================================================
\section*{R-4: Log-cancellation proof wrong}

\rev{Reviewer: The winding-number argument at L1115--1118 is mathematically
false. $F = -q$ lies inside the unit circle, giving winding number 1.}

\textbf{Response:} The referee is correct. The winding-number statement was
incorrect. We have replaced it with a rigorous residue-theorem argument:
\begin{enumerate}[nosep]
\item The zero-mean-subtracted integrand has vanishing contour integral
  by construction.
\item By the residue theorem, the sum of residues inside $|F|=1$ (at
  $F=0$ and $F=-q$) must equal zero.
\item The residue at $F=0$ vanishes because the $M$-average has been
  subtracted.
\item Therefore the residue at $F=-q$ also vanishes.
\item The $\log(1+qF)$ coefficient vanishes independently by the
  periodicity requirement.
\end{enumerate}
A structural remark has been added noting that this vanishing is a
consequence of the zero-mean normalization, not a coincidence.

\textbf{Changes:} Section~8.1 (antiderivative assembly): proof completely
rewritten with contour-integral equation and residue-theorem argument.

%% ====================================================================
\section*{R-5: $w_h$ formula missing $\gamma$}

\rev{Reviewer: The display formula $w_h = (3A/cr^3)\,c_i\,\hat z^2$ is
missing a factor of $\gamma$. The correct expression is
$w_h = 3A\delta\,\hat z^2/(cr^3)$.}

\textbf{Response:} The referee correctly identified a missing factor of
$\gamma$ in the display formula (since $\delta = \gamma c_i$). The formula
has been corrected, and a derivation sketch has been added showing the chain
from the exact GEqOE definitions to the first-order result.
The reduced equations and all numerical results are unaffected because
the symbolic computation operates on the exact GEqOE equations, not on the
hand-reduced display formula.

\textbf{Changes:} Eq.~(wh\_j2\_first): corrected $c_i \to \delta$.
New derivation paragraph with explicit intermediate steps.

%% ====================================================================
\section*{R-6: Paper not self-contained}

\rev{Reviewer: Placeholders and deferrals to companion notes make the paper
incomplete.}

\textbf{Response:} We agree. Four specific improvements:
\begin{enumerate}[nosep]
\item The $\langle F^k \rangle_M$ placeholder has been replaced with
  the explicit contour-integral formula and the first two values
  ($\langle F \rangle_M = -g$, $\langle F^2 \rangle_M = g^2(1+2\beta)/(1+\beta)^2$).
\item The J$_3$/J$_4$ drift formulas are too long for inline display
  (500--1000+ character rational expressions), but the paper now states
  their harmonic structure and characterizes the coefficient functions
  as rational functions of degree 8--16 in $(q, Q)$.
\item The general degree-$n$ result is now summarized in the paper
  (structure of the residue computation, linear growth of degrees).
\item All companion-note references have been changed to ``supplementary
  material'' or ``electronic appendix.''
\end{enumerate}

\textbf{Changes:} Eq.~(mean\_Fk): explicit formulas. Section~10.4:
expanded coefficient characterization. Section~10.4: general degree-$n$
structural result. All \texttt{docs/geqoe\_averaged/} paths removed.

%% ====================================================================
\section*{R-7: Performance ranges numerically inconsistent}

\rev{Reviewer: Stated improvement ranges (4--80$\times$ for Brouwer,
1.5--8$\times$ for DSST) do not match the tables.}

\textbf{Response:} The referee is correct. We have recomputed all
improvement ratios from the published tables:
\begin{itemize}[nosep]
\item GEqOE vs Brouwer: 4--335$\times$ (was 4--80$\times$).
\item GEqOE vs DSST: 1.5--93$\times$ (was 1.5--8$\times$).
\end{itemize}
All instances in the abstract, conclusions, and appendix discussion have
been corrected. The ``best case 335$\times$'' is consistent with
Table~\ref{tab:dsst\_osculating} ($2.884/0.0086 = 335$).

\textbf{Changes:} Abstract, Section~13 (conclusions), Appendix~A
discussion, Appendix~B discussion: all ratios corrected.

%% ====================================================================
\section*{R-8: Novelty positioning too loose}

\rev{Reviewer: The introduction does not isolate the real novelties
sharply enough.}

\textbf{Response:} The novelty paragraph has been rewritten to state three
specific respects in which the present work differs from prior art:
(1)~direct operation in non-canonical GEqOE variables,
(2)~averaging in the generalized eccentric longitude $K$,
(3)~short-period construction by residue decomposition in the complex true
anomaly $F$.

\textbf{Changes:} Section~1 (Introduction): novelty paragraph completely
rewritten.

%% ====================================================================
\section*{R-9: $w$ notation reuse}

\rev{Reviewer: $w = \sqrt{\mu/a}$ and $w = e^{i\omega}$ create ambiguity.}

\textbf{Response:} The complex pericenter exponential has been renamed to
$W = e^{i\omega}$ (capital $W$) throughout Sections~7--8 and the residue
decomposition. The mean-motion quantity $w = \sqrt{\mu/a}$ is unchanged.

\textbf{Changes:} Eq.~(W\_def) and all occurrences in Sections~7--8:
$w \to W$.

%% ====================================================================
\section*{R-10: Internal-note language and repo paths}

\rev{Reviewer: Several passages use internal-note language and repository
paths.}

\textbf{Response:} All instances have been fixed:
\begin{itemize}[nosep]
\item ``current implementation'' $\to$ removed
\item ``current branch uses'' $\to$ ``the exact conservative GEqOE system is''
\item ``present branch'' $\to$ ``present variables''
\item ``current probe'' $\to$ ``Numerically''
\item All \texttt{docs/geqoe\_averaged/} paths $\to$ ``supplementary material''
\item ``This note'' $\to$ ``This paper'' throughout
\end{itemize}

\textbf{Changes:} Sections~2, 3, 10, and throughout.

%% ====================================================================
\section*{R-11: Phipps citation}

\rev{Reviewer: ``Warren Phipps' 1992 fix'' is not a proper citation.}

\textbf{Response:} Replaced with reference to Orekit's internal
critical-inclination correction as documented in the Orekit v13.1 API
documentation.

\textbf{Changes:} Appendix~A (Brouwer comparison): Phipps reference
replaced with Orekit citation.

%% ====================================================================
\section*{R-12: $q$ notation overloading}

\rev{Reviewer: The eccentricity parameter $q$ could be confused with the
inclination components $q_1, q_2$.}

\textbf{Response:} A parenthetical disambiguation has been added after the
definition of $q = g/(1+\beta)$.

\textbf{Changes:} Eq.~(q\_def): parenthetical remark added.

%% ====================================================================
\section*{R-13: $\omega$ osculating vs mean not flagged}

\rev{Reviewer: $\omega = \Psi - \Omega$ is used interchangeably for
osculating and mean values without clarification.}

\textbf{Response:} A clarifying statement has been added at the beginning of
Section~6.3: ``In the integrands that follow, all slow variables---including
$\omega$, $g$, $Q$---are evaluated at their mean values, consistent with the
first-order averaging prescription.''

\textbf{Changes:} Section~6.3 opening: new paragraph.

%% ====================================================================
\section*{R-14: Sign/definition chain not shown}

\rev{Reviewer: The chain $A \to \mathcal{E}_Z \to \dot{p}$ is not shown
explicitly.}

\textbf{Response:} A derivation paragraph has been added after the corrected
$w_h$ formula, showing how the out-of-plane forcing $F_h$ reduces to the
first-order expression via $\hat z$ and $\delta = \gamma c_i$.

\textbf{Changes:} Section~6.3: new derivation paragraph after
Eq.~(wh\_j2\_first).

%% ====================================================================
\section*{R-15: $\dot\Psi$ narrative confusion}

\rev{Reviewer: The cautionary remark about equations (86)/(88) is confusing.}

\textbf{Response:} A confirmation statement has been added: ``Equations~(86)
and the above are algebraically consistent, confirming that the direct
derivation from the $\dot p$ equations reproduces the expected structure.''

\textbf{Changes:} After Eq.~(Psidot\_j2\_first) derivation: one sentence
added.

%% ====================================================================
\section*{R-16: Averaging weight $T = 2\pi a$ unexplained}

\rev{Reviewer: Why is $T = 2\pi a$ in the $K$-average?}

\textbf{Response:} A derivation paragraph has been added explaining that
$T = \int_0^{2\pi}(r/w)\,dK = 2\pi a$ follows from $\langle r \rangle_K = a$
for the unperturbed ellipse, and noting that the weight $r/a$ replaces the
unit weight of the classical $M$-average.

\textbf{Changes:} After Eq.~(avg\_operator): new paragraph.

%% ====================================================================
\section*{R-17: Missing general-zonal $\dot{\bar M}$}

\rev{Reviewer: The general-zonal mean anomaly rate is absent from the paper
(though implemented in code).}

\textbf{Response:} The referee correctly identified a gap in the paper.
Eq.~(Mdot\_total\_N) has been added after the mean-state ODE, giving
$\dot{\bar M}$ in the same finite-harmonic form as the slow variables.
The runtime procedure (Step~2) now explicitly states that $\bar M$ is
integrated alongside the slow state. Step~3 clarifies the Kepler equation
solve for $\bar K$ from $\bar M$.

\textbf{Changes:} New Eq.~(Mdot\_total\_N) with surrounding text.
Section~9.2, Steps 2--3: updated.

%% ====================================================================
\section*{R-18: $J_5$ claim unsupported}

\rev{Reviewer: The claim that results are ``indistinguishable from
$J_2$--$J_4$'' is unsupported by table data.}

\textbf{Response:} The text now quotes the maximum difference numerically:
``The maximum position RMS difference between $J_2$--$J_4$ and $J_2$--$J_5$
across all test cases is less than 1~meter.''

\textbf{Changes:} Section~11.3: quantified the $J_5$ contribution.

%% ====================================================================
\section*{R-19: $\ell_1$ derivation gap}

\rev{Reviewer: The derivation of Eq.~(ell\_v\_relation) is not shown.}

\textbf{Response:} A one-line remark has been added: ``Equation~(ell\_v\_relation)
follows from expanding $L_{\mathrm{osc}}$ at first order using
$\partial L/\partial K = r/a$; the cross-term
$v_1(p_1\sin K + p_2\cos K)$ reduces to $(\bar r/a)\,v_1$.''

\textbf{Changes:} After Eq.~(ell\_v\_relation): one-line remark.

%% ====================================================================
\section*{R-20: $n_0^2$ division well-definedness}

\rev{Reviewer: The division by $n_0 = w/r$ in the $v_1$ quadrature may
not be well-defined.}

\textbf{Response:} A parenthetical note has been added: ``(the division
by $n_0 = w/r$ is well-defined for all non-rectilinear orbits, where
$r > 0$ throughout).''

\textbf{Changes:} After Eq.~(v\_prime): parenthetical note.

%% ====================================================================
\section*{R-21: Brouwer comparison fairness}

\rev{Reviewer: Since Brouwer is second-order in $J_2$, the 4--335$\times$
advantage may be from better SP reconstruction, not secular drift.}

\textbf{Response:} A comparison fairness paragraph has been added in the
Appendix~A discussion, noting that Brouwer--Lyddane is a purely analytical
theory with no mean-only output mode, so the reported errors unavoidably
combine secular drift and short-period reconstruction. The DSST mean-drift
decomposition (Table~8) provides indirect evidence that short-period
reconstruction is the dominant error source, though this cannot be
separated for Brouwer directly.

\textbf{Changes:} Appendix~A, Discussion: new ``Comparison fairness''
paragraph.

%% ====================================================================
\section*{R-22: $h$ domain of validity}

\rev{Reviewer: When is the radicand in $h = \sqrt{c^2 - 2r^2 U}$ positive?}

\textbf{Response:} A remark has been added: ``The radicand is positive for
perturbations satisfying $2r^2|U|/c^2 \ll 1$, which holds for all practical
zonal truncations.''

\textbf{Changes:} After the exact GEqOE system: one-line remark.

%% ====================================================================
\section*{R-23: Thrust section too thin}

\rev{Reviewer: Section~12 is a brief sketch with no validation, weakening
the otherwise rigorous presentation.}

\textbf{Response:} We agree. The standalone thrust section has been removed.
The key structural observation has been compressed into a paragraph in the
Conclusions, framed as future work: the GEqOE structure makes thrust-induced
averaged kernels low-order trigonometric polynomials in $K$.

\textbf{Changes:} Section~12 removed. Section~13 (Conclusions): thrust
extension described in extension item~(iii).

%% ====================================================================
\section*{R-24: Overclaiming language}

\rev{Reviewer: ``Complete,'' ``exact,'' and ``orders of magnitude'' are used
without sufficient qualification.}

\textbf{Response:} ``Orders of magnitude'' claims that were not backed by
specific numbers have been qualified with references to the relevant
tables or with implementation caveats. ``This note'' has been changed to
``This paper'' throughout. The timing claim in the DSST appendix has been
qualified with the measured wall-clock ratio.

\textbf{Changes:} Abstract, Appendix~B discussion and caveat: language
audited and qualified.

%% ====================================================================
\section*{R-25: $\backslash$date\{$\backslash$today\} on title page}

\rev{Reviewer: The date should not be \texttt{$\backslash$today}.}

\textbf{Response:} Replaced with ``March 2026.''

\textbf{Changes:} Line~21: fixed date.

%% ====================================================================
\section*{R-26: Repetitive phrasing}

\rev{Reviewer: ``No numerical quadrature, no interpolation, no iterative
solution'' appears multiple times.}

\textbf{Response:} The phrase now appears only twice: once in the abstract
and once in Section~9 (operational structure), which is acceptable.

\textbf{Changes:} Verified; no additional removals needed.

%% ====================================================================
\section*{R-27: Missing limitations paragraph}

\rev{Reviewer: The paper does not discuss its limitations.}

\textbf{Response:} A limitations paragraph has been added before the
extensions list in the Conclusions, covering:
(1)~first-order truncation ($O(\varepsilon^2)$ neglected),
(2)~singular cases (retrograde equatorial, rectilinear),
(3)~high-eccentricity degradation,
(4)~zonal-only restriction.

\textbf{Changes:} Section~13 (Conclusions): new ``Limitations'' paragraph
with 4-item list.

%% ====================================================================
\section*{R-28: Timing claim misleading}

\rev{Reviewer: ``10--100$\times$ faster'' is misleading when the Python
prototype is actually slower than compiled DSST.}

\textbf{Response:} The timing claim has been qualified throughout:
\begin{itemize}[nosep]
\item Abstract: now states ``theoretical per-epoch cost that is orders of
  magnitude smaller (though the current Python prototype is slower than the
  compiled Java DSST).''
\item Conclusions: ``theoretical per-epoch cost that is orders of magnitude
  smaller (Table, implementation caveat).''
\item Appendix~B caveat: now states the measured wall-clock ratio explicitly
  (``6--300$\times$ slower than compiled DSST'').
\end{itemize}

\textbf{Changes:} Abstract, Section~13, Appendix~B: timing language qualified.

\end{document}
